% The diagram is the actual finite tensor contraction on a rooted tree.
\begin{tikzpicture}[x=1cm,y=1cm,line cap=round,line join=round,
  font=\small,
  vertex/.style={circle,draw=blue!65!black,fill=blue!8,
    minimum size=8mm,inner sep=1pt},
  edge/.style={line width=1pt,draw=blue!65!black},
  sub/.style={draw=gray!65,fill=gray!7,line width=.7pt}]
\node[vertex] (parent) at (0,3.15) {$u$};
\node[above=5pt] at (parent.north) {Parent angle};
\node[vertex] (v) at (0,.8) {$t$};
\draw[edge] (parent)--(v);
\node[fill=white,inner sep=3pt,align=center] at (-.9,2.0)
 {$K_N(t,u;s)$};
\node[right=8pt,align=left] at (v.east)
 {Vertex $v$\\crossing power $n_v$};
\coordinate (left) at (-1.8,-.65);
\coordinate (middle) at (0,-.65);
\coordinate (right) at (1.8,-.65);
\foreach \p in {left,middle,right}{
  \draw[edge] (v)--(\p);
  \path[sub] (\p)--($(\p)+(-.55,-1.25)$)--
    ($(\p)+(.55,-1.25)$)--cycle;
}
\node at (-1.8,-1.5) {$w_1$};
\node at (0,-1.5) {$\cdots$};
\node at (1.8,-1.5) {$w_m$};
\node[fill=white,inner sep=2pt,rotate=39] at (-1.3,-.15)
 {$M_{w_1,N}(t)$};
\node[fill=white,inner sep=2pt,rotate=-39] at (1.3,-.15)
 {$M_{w_m,N}(t)$};
\node[align=center,text width=5.2cm] at (0,-3.05)
 {Each child subtree has already been evaluated\\
 as a polynomial in $-2\cos t$.};
\node[anchor=north west,align=left,text width=7.4cm] (kerneltext)
 at (3.3,3.65) {
 The edge keeps every index $0\le k<N$:
 \[
 K_N(t,u;s)=\sum_{k=0}^{N-1}
 e_k(-2\cos t)e_k(-2\cos u)\frac{H_N}{W_k}.
 \]
 The trace at $v$ acts only in $t$:
 \[
 M_{v,N}(u)=\gamma_{e_v}\,
 \mathcal D_v^{(t)}
 \left[K_N(t,u;s)\prod_{j=1}^{m}M_{w_j,N}(t)\right].
 \]
 Trigonometric degree in $u$: at most $N-1$.
 };
\node[anchor=north west,align=left,text width=7.4cm]
 at ([yshift=-4mm]kerneltext.south west) {
 Factor out the exact value at zero:
 \[
 \widehat M_{v,N}(u)=\frac{M_{v,N}(u)}{M_{v,N}(0)}.
 \]
 Its local phase and its real decay control the next
 trace. The denominator includes all crossing and
 orientation factors from this subtree.
 };
\end{tikzpicture}
